% Suggested replacement text for the TWO additions reported by the author.
% These fragments have NOT been inserted into the source manuscript.
% Requires xcolor, natbib, and the accompanying AE1_references.bib.

% Section 1.2 Related Works: place at the end of this subsection,
% immediately before Section 2 Dataset Compilation.
{\color{blue}
\noindent\textbf{Knowledge and Skills in LLM Agents.}
Prior work distinguishes declarative knowledge about facts and
concepts from procedural knowledge about how to perform a task
\citep{ae1_li2024metacognitive}. Retrieval-augmented generation
conditions generation on external material
\citep{ae1_lewis2020rag}; the retrieved content can include
input--output examples \citep{ae1_rubin2022epr} or reusable
reasoning templates \citep{ae1_yang2024buffer}. Retrieval therefore
describes an information-access mechanism rather than a particular
type of knowledge. Recent research characterizes agent skills as
reusable procedural artifacts with instructions, supporting
resources, and applicability conditions
\citep{ae1_zhou2026skillsurvey}. Skill packages can contain
reference material and executable components
\citep{ae1_li2026skillsbench,ae1_openai2026skills}, and can guide
retrieval during execution \citep{ae1_openai2026skilltools}.
Knowledge and skills are consequently complementary: reference
content can support a reusable procedure, while retrieval can supply
the information needed to execute it. Agent Workflow Memory induces
workflows from task experience \citep{ae1_wang2025awm}, and
OptSkills learns and refines optimization procedures organized by
problem archetypes \citep{ae1_yang2026optskills}. LEAN-LLM-OPT
shares the emphasis on procedural reuse, while obtaining its
reference resources and reusable modeling guidance through
design-time curation.
}

% Section 3.3 Novelties in Our Design and Key Messages:
% place immediately after the paragraph about harnessing LLM flexibility
% through structured workflows, and before the next existing key message.
{\color{blue}
\noindent\textbf{Connecting Reference Knowledge and Modeling Skills.}
LEAN-LLM-OPT combines case-based modeling knowledge with reusable
procedures for accessing data and constructing formulations.
Ref-Data supplies reference descriptions, sample data, expert
formulations, and annotations of relevant data. Reference retrieval
supports classification and demonstration construction, whereas
CSVQA obtains information from the target datasets for the current
formulation. In the type-tailored route, the workflow generation
agent combines selected reference content with modeling instructions
and tool-use steps to construct a demonstration. The model generation
agent adapts this guidance to the target requirements and data.
These procedures are skill-like in their reusable organization of
modeling knowledge and data-access operations. Their reuse does not
require identical execution across instances: the reference bank
and reusable instructions remain fixed during evaluation, while
the selected reference and subsequent reasoning and tool use can
vary. For Others and Mixture, the type-agnostic route uses an
instance-specific abstract model plan. LEAN-LLM-OPT implements
procedural reuse through prompts, workflow-construction code, and
data tools; it does not package these procedures as standalone
skills or learn a persistent skill library from completed
trajectories. Its methodological focus is the coordinated use of
reference knowledge and procedural guidance to formulate optimization
models from problem descriptions and external instance data.
}
